\setcounter{secnumdepth}{3}
\titlespacing*{\chapter}{0pt}{-0.1cm}{1cm}
\titleformat{\chapter}[display]
  {\bfseries\LARGE}  
  {Phần~\thechapter}{0pt}{\raggedright}

\chapter{Overfitting: Bản chất và cách nhận diện}

\section{Giới thiệu về Overfitting}

Trong lĩnh vực Machine Learning, đặc biệt là học có giám sát (supervised learning), một trong những vấn đề quan trọng ảnh hưởng trực tiếp đến hiệu quả của mô hình là hiện tượng \textit{overfitting}. Overfitting xảy ra khi mô hình học quá chi tiết từ dữ liệu huấn luyện, bao gồm cả các nhiễu (noise), dẫn đến khả năng tổng quát hóa kém trên dữ liệu chưa từng thấy. Nói cách khác, mô hình có thể đạt độ chính xác rất cao trên tập huấn luyện nhưng lại hoạt động kém trên tập kiểm tra \cite{Ying2019}.

Về bản chất, mục tiêu của Machine Learning không phải là tối ưu hóa tuyệt đối trên dữ liệu huấn luyện, mà là tối đa hóa khả năng dự đoán trên dữ liệu mới. Tuy nhiên, nếu mô hình được tối ưu hóa quá mức để giảm sai số trên tập huấn luyện, nó có xu hướng ghi nhớ các đặc điểm riêng lẻ hoặc nhiễu của dữ liệu thay vì học được quy luật tổng quát. Điều này dẫn đến sự sai lệch giữa mục tiêu tối ưu hóa và mục tiêu thực tế của bài toán học máy \cite{Dietterich1995}.

Tầm quan trọng của việc xử lý overfitting nằm ở khả năng quyết định tính ứng dụng thực tế của mô hình. Một mô hình chỉ hoạt động tốt trên dữ liệu huấn luyện nhưng thất bại trên dữ liệu thực tế sẽ không có giá trị sử dụng. Do đó, việc nhận diện và kiểm soát overfitting là một kỹ năng cốt lõi trong quá trình xây dựng mô hình học máy.

\section{Bản chất của Overfitting}

Overfitting là một khái niệm mang tính tương đối và so sánh. Một mô hình được coi là bị overfit khi nó có hiệu suất tốt hơn trên tập huấn luyện nhưng lại kém hơn trên dữ liệu ngoài mẫu so với một mô hình khác, hoặc so với chính nó tại một thời điểm trước đó trong quá trình huấn luyện.

Về bản chất, overfitting phản ánh sự mất cân bằng giữa việc học dữ liệu huấn luyện và khả năng tổng quát hóa. Khi mô hình quá tập trung vào việc giảm sai số huấn luyện, nó có xu hướng ghi nhớ dữ liệu thay vì học quy luật tổng quát.

Cần phân biệt rõ giữa overfitting và khả năng tổng quát hóa kém. Một mô hình có thể có sai số cao trên cả $E_{in}$ và $E_{out}$, nhưng trường hợp này không được xem là overfitting. Overfitting chỉ xảy ra khi mô hình tối ưu hóa quá mức $E_{in}$ dẫn đến làm giảm hiệu suất trên $E_{out}$.

\section{Dấu hiệu nhận diện Overfitting}

Một trong những dấu hiệu quan trọng nhất để nhận diện overfitting là mối quan hệ giữa sai số huấn luyện và sai số kiểm tra. Theo Abu-Mostafa trong \textit{Learning From Data} (Lecture 11: \textit{Overfitting}), khi sai số trong mẫu \(E_{in}\) tiếp tục giảm nhưng sai số ngoài mẫu \(E_{out}\) lại tăng (\textit{``Ein $\downarrow$ Eout $\uparrow$''}), đó là dấu hiệu mô hình đang rơi vào hiện tượng overfitting \cite{caltech_lfd}.

Hiện tượng này thường xuất hiện rõ ràng trong quá trình huấn luyện các mô hình học sâu. Ban đầu, cả $E_{in}$ và $E_{out}$ đều giảm do mô hình dần học được các đặc trưng quan trọng. Tuy nhiên, sau một số epoch nhất định, việc tiếp tục huấn luyện sẽ khiến mô hình bắt đầu ghi nhớ dữ liệu huấn luyện, dẫn đến việc $E_{out}$ tăng lên trong khi $E_{in}$ vẫn tiếp tục giảm.

Ngoài ra, một số biểu hiện khác của overfitting có thể bao gồm:
\begin{itemize}
    \item Mô hình quá phức tạp so với lượng dữ liệu.
    \item Hiệu suất trên dữ liệu mới kém dù kết quả trên tập huấn luyện rất tốt.
\end{itemize}

\section{Vai trò của nhiễu trong Overfitting}

Một nguyên nhân cốt lõi của overfitting là việc mô hình học cả nhiễu trong dữ liệu. Theo \cite{Ying2019}, các mô hình bị overfit có xu hướng ghi nhớ toàn bộ dữ liệu huấn luyện, bao gồm cả những thành phần không mang tính quy luật.

Nhiễu trong dữ liệu có thể đến từ nhiều nguồn khác nhau, chẳng hạn như sai số đo lường, dữ liệu không đầy đủ hoặc các biến ngẫu nhiên không thể dự đoán. Khi mô hình cố gắng khớp với các dữ liệu này, nó sẽ học những mẫu không có ý nghĩa thực tế, từ đó làm giảm hiệu suất tổng thể.

Do đó, việc kiểm soát mức độ phức tạp của mô hình và hiểu rõ vai trò của nhiễu là yếu tố quan trọng trong việc hạn chế overfitting.

\section{Tổng hợp các dấu hiệu nhận biết}

Từ các phân tích trên, có thể tổng hợp một số dấu hiệu chính để nhận diện overfitting như sau:
\begin{itemize}
    \item Sai số huấn luyện giảm nhưng sai số kiểm tra tăng.
    \item Mô hình quá phức tạp so với dữ liệu.
    \item Mô hình học cả nhiễu thay vì quy luật.
    \item Hiệu suất trên dữ liệu mới kém dù kết quả trên dữ liệu huấn luyện rất tốt.
\end{itemize}